% Part: incompleteness
% Chapter: theories-computability
% Section: computably-axiomatizable

\documentclass[../../../include/open-logic-section]{subfiles}

\begin{document}

\olfileid{inc}{tcp}{cax}

\olsection{النظريات \printtoken{S}{axiomatizable}}

تكون النظرية $\Th{T}$ \emph{!!{axiomatizable}} إذا أمكن اختيار
مجموعة بديهيات قابلة للحساب $\OLId{latin-looped}{A}$ لها. ومعنى أن $\OLId{latin-looped}{A}$ مجموعة بديهيات
لـ$\Th{T}$ هو $\OLId{latin-looped}{T}=\Setabs{!A}{\OLId{latin-looped}{A}\Proves !A}$. وهذه الخاصية تتحقق في
كل أكسمة «معقولة» للأعداد الطبيعية؛ وبخاصة، كل نظرية يكفيها عدد
منته من البديهيات تكون !!{axiomatizable}.

\begin{lem}
إذا كانت $\Th{T}$ !!{axiomatizable}، لزم أن تكون $\Th{T}$
!!{computably
enumerable}.
\end{lem}

\begin{proof}
لتكن $\OLId{latin-looped}{A}$ مجموعة بديهيات قابلة للحساب للنظرية $\Th{T}$. وشبه القرار
في مسألة $!A\in \OLId{latin-looped}{T}$ يحصل بالبحث عن !!a{derivation} للصيغة $!A$
من تلك البديهيات.

وبيان ذلك على وجه آخر: انتماء $!A$ إلى $\Th{T}$ يكافئ وجود قائمة
منتهية من البديهيات $!B_\OLMathNumeral{1}$، \dots، $!B_\OLId{latin-ordinary}{k}$ من $\OLId{latin-looped}{A}$، مع
!!a{derivation} للصيغة
$(!B_\OLMathNumeral{1}\land\dots\land !B_\OLId{latin-ordinary}{k})\lif !A$ في منطق الرتبة الأولى.
وقد تقدم أن المجموعة المعرّفة بقولنا «يوجد \dots بحيث \dots»
تكون !!{c.e.} متى كان الشرط الذي تشغل موضعه «\dots» الثانية قابلًا
للحساب.
\end{proof}

\end{document}
